% Minimal preamble for Open Logic Project sources (pinned commit 1e960beff9ed7835bf3e3f1335e21af3439cd107).
%
% OLP's own master preamble cannot run on the project host: TeX Live 2025 there lacks
% prooftree.sty, ntheorem.sty, extarrows.sty, mathdots.sty, isomath.sty and subfiles.sty.
% Following the project's own rule ("if the author preamble needs packages the toolchain lacks,
% do not invent macros -- use a minimal working preamble"), this file declares, verbatim in
% spirit and numbering behaviour, exactly the theorem-like environments OLP defines in
% sty/open-logic-book-envs.sty (thm/lem/prop/cor/defn/ex/prob/rem/note/case/axiom/conv),
% plus the shorthand accents OLP uses (`!!a`, `!!b`, `!!{...}`).
%
% Author text is untouched; only the packaging is ours.
\documentclass[10pt,a4paper]{article}

\usepackage{amsmath}
\usepackage{amssymb}
\usepackage{amsthm}
\usepackage{graphicx}
\usepackage{hyperref}
\usepackage{enumitem}
\usepackage[framemethod=default]{mdframed}

% ---- OLP shorthand accents (from sty/open-logic.sty) ----
% NB: \! is taken by LaTeX (negative thin space), so we only define the !!a / !!b families
% and the \olup helper used by them.
\newcommand*{\olup}[1]{\MakeUppercase{#1}}
\def\!!!a#1{\olup{#1}}
\def\!!a#1{\olup{#1}}
\def\!!b#1{\olup{#1}}

% ---- theorem-like environments, same family numbering as OLP ----
\theoremstyle{plain}
\newtheorem{thm}{Theorem}[section]
\newtheorem{lem}[thm]{Lemma}
\newtheorem{prop}[thm]{Proposition}
\newtheorem{cor}[thm]{Corollary}
\newtheorem{axiom}{Axiom}

\theoremstyle{definition}
\newtheorem{defn}[thm]{Definition}
\newtheorem{ex}[thm]{Example}
\newtheorem{prob}{Problem}[section]
\newtheorem{probd}{Problem}
\newtheorem{conv}[thm]{Convention}

\theoremstyle{remark}
\newtheorem{rem}[thm]{Remark}
\newtheorem{note}[thm]{Note}
\newtheorem{case}[thm]{Case}

% ---- environments OLP defines as wrappers ----
\newenvironment{probtag}[1]{\begin{prob}}{\end{prob}}
\newenvironment{explain}{}{}
\newenvironment{explain*}{}{}
\newenvironment{editorial}{}{}
\newenvironment{digress}{}{}
\newenvironment{digress*}{}{}
\newenvironment{history}{}{}
\newenvironment{reference}{}{}
\newcommand*{\olfileid}[3]{}
\newcommand*{\olsection}[1]{\section*{#1}}
\newcommand*{\olchapter}[3]{}
\newcommand*{\olpart}[2]{}
\newcommand*{\olimport}[1]{}
\newcommand*{\olinclude}[1]{}
\newcommand*{\olpath}[1]{}
\newcommand*{\ollangid}[1]{}
\newcommand*{\problemtag}[1]{}
\newcommand*{\problemauthor}[1]{}
\newcommand*{\ollabel}[1]{\label{#1}}
% \olref takes 0-3 OPTIONAL arguments plus one mandatory label:
%     \olref{lab}  |  \olref[from]{lab}  |  \olref[from][to]{lab}  |  \olref[a][b][c]{lab}
% A single-argument definition silently swallows "[sth]" and emits a broken reference
% (measured on the host: the text came out as "??iso]x"). The chain below consumes the
% brackets and cross-references only the mandatory argument.
\makeatletter
\newcommand{\olref}{\@ifnextchar[{\olref@i}{\olref@final}}
\def\olref@i[#1]{\@ifnextchar[{\olref@ii}{\olref@final}}
\def\olref@ii[#1]{\@ifnextchar[{\olref@iii}{\olref@final}}
\def\olref@iii[#1]{\olref@final}
\def\olref@final#1{\ref{#1}}
\makeatother
\newcommand*{\olnewpage}{\clearpage}

% ===================== OLP compatibility layer (1) =====================

\makeatletter

% ---------------- proof trees (OLP's prooftree is a bussproofs derivative) ----------------
\providecommand{\AxiomC}[1]{}
\providecommand{\UnaryInfC}[1]{}
\providecommand{\BinaryInfC}[1]{}
\providecommand{\TrinaryInfC}[1]{}
\providecommand{\QuaternaryInfC}[1]{}
\providecommand{\QuinaryInfC}[1]{}
\providecommand{\RightLabel}[1]{}
\providecommand{\LeftLabel}[1]{}
\providecommand{\doubleLine}{}
\providecommand{\noLine}{}
\providecommand{\DisplayProof}{}
\providecommand{\fCenter}{}
\newenvironment{prooftree}{}{}
\newenvironment{ltprooftree}{}{}
\newenvironment{oltableau}{}{}

% ---------------- logical notation ----------------
\providecommand{\lif}{\mathbin{\rightarrow}}
\providecommand{\lforall}{\mathord{\forall}}
\providecommand{\lexists}{\mathord{\exists}}
\providecommand{\ltrue}{\top}
\providecommand{\lfalse}{\bot}
\providecommand{\Ltrue}{\top}
\providecommand{\Lfalse}{\bot}
\providecommand{\Sequent}{\mathrel{\Rightarrow}}
\providecommand{\Proves}{\mathrel{\vdash}}
\providecommand{\Entails}{\mathrel{\models}}
\providecommand{\Subst}[3]{#1[#2/#3]}
\providecommand{\sFmla}[2]{#1(#2)}
\providecommand{\pFmla}[2]{#1(#2)}
\providecommand{\formula}{\relax}
\providecommand{\True}{\top}
\providecommand{\False}{\bot}
\providecommand{\Args}[1]{#1}
\providecommand{\Ts}{\ensuremath{\top}}
\providecommand{\Fi}{\ensuremath{\bot}}
\providecommand{\lfalsum}{\bot}

% ---------------- objects, structures, metatheory ----------------
\providecommand{\Obj}{}
\providecommand{\Struct}{\ensuremath{\mathfrak{A}}}
\providecommand{\Domain}{\ensuremath{|\mathfrak{A}|}}
\providecommand{\Assign}{\ensuremath{s}}
\providecommand{\Sat}[2]{\ensuremath{#1 \models #2}}
\providecommand{\mSat}[3]{\ensuremath{#1 \models #2[#3]}}
\providecommand{\mModel}[2]{\ensuremath{#1 \models #2}}
\providecommand{\Lang}{\ensuremath{\mathcal{L}}}
\providecommand{\Th}{\ensuremath{\mathrm{Th}}}
\providecommand{\Log}{\ensuremath{\mathrm{L}}}
\providecommand{\tuple}[1]{\langle #1 \rangle}
\providecommand{\num}[1]{\underline{#1}}
\providecommand{\cfind}[1]{\varphi_{#1}}
\providecommand{\Nat}{\ensuremath{\mathbb{N}}}
\providecommand{\Mult}{\ensuremath{\mathrm{Mult}}}
\providecommand{\Add}{\ensuremath{\mathrm{Add}}}
\providecommand{\Exp}{\ensuremath{\mathrm{Exp}}}
\providecommand{\Zero}{\ensuremath{0}}
\providecommand{\Succ}{\ensuremath{\mathrm{S}}}
\providecommand{\lambd}{\ensuremath{\lambda}}
\providecommand{\bredpar}{\mathrel{\rightarrow_\beta}}
\providecommand{\bcd}[1]{\ensuremath{\ulcorner #1 \urcorner}}
\providecommand{\Set}[1]{\ensuremath{\{#1\}}}
\providecommand{\indcase}[1]{\par\noindent\emph{Case #1}.\quad}
\providecommand{\indfrm}[1]{\par\noindent\emph{Form #1}.\quad}
\providecommand{\iftag}[1]{#1}
\providecommand{\tagitem}[1]{\item[#1]}
\providecommand{\Elim}[1]{}
\providecommand{\Intro}[1]{}
\providecommand{\LeftR}[1]{}
\providecommand{\RightR}[1]{}
\providecommand{\TAss}[1]{}
\providecommand{\TRule}[1]{}

\makeatother
\providecommand{\olpath}{.}

% ===================== OLP compatibility layer (2) =====================
% --- second compatibility block: notation and shorthands used by the frozen OLP units ---
\providecommand{\CutCS}{\mathrm{Cut}}
\providecommand{\LeftSubproofLabel}{\relax}
\providecommand{\RightSubproofLabel}{\relax}
\providecommand{\maeh}{\relax}
\providecommand{\cheight}{\relax}
\providecommand{\ande}{\mathrm{and}_E}
\providecommand{\andi}{\mathrm{and}_I}
\providecommand{\ore}{\mathrm{or}_E}
\providecommand{\ori}{\mathrm{or}_I}
\providecommand{\typeof}{\mathord{:}}
\providecommand{\pValue}[1]{#1}
\providecommand{\Sent}{\mathrm{Sent}}
\providecommand{\cardplus}{\mathbin{+}}
\providecommand{\cardtimes}{\mathbin{\cdot}}
\providecommand{\cardexpo}{\mathbin{\hat{\ }}}
\providecommand{\cardnless}{\not<}
\providecommand{\cardle}{\leq}
\providecommand{\funfromto}[3]{#1\colon #2 \to #3}
\providecommand{\funimage}[1]{\operatorname{im}(#1)}
\providecommand{\fdefined}{\mathord{\downarrow}}
\providecommand{\canonord}{\operatorname{ord}}
\providecommand{\ordtimes}{\mathbin{\cdot}}
\providecommand{\ordeq}[2]{#1 = #2}
\providecommand{\ordneq}[2]{#1 \neq #2}
\providecommand{\ordtype}[1]{\operatorname{ord}(#1)}
\providecommand{\ordsucc}{\mathord{+1}}
\providecommand{\ordunion}{\cup}
\providecommand{\onesphere}{S^1}
\providecommand{\rotationsgroup}{\mathrm{SO}(2)}
\providecommand{\equivrep}[1]{[#1]}
\providecommand{\rlexless}{\prec}
\providecommand{\ZFC}{\mathrm{ZFC}}
\providecommand{\ZF}{\mathrm{ZF}}
\providecommand{\Rat}{\mathbb{Q}}
\providecommand{\Real}{\mathbb{R}}
\providecommand{\PosInt}{\mathbb{Z}^{+}}
\providecommand{\Id}{\mathrm{id}}
\providecommand{\Bin}{\mathrm{Bin}}
\providecommand{\pto}{\rightharpoonup}
\providecommand{\dom}{\operatorname{dom}}
\providecommand{\ran}{\operatorname{ran}}
\providecommand{\len}{\operatorname{len}}
\providecommand{\concat}{\mathbin{{}^\frown}}
\providecommand{\card}[1]{\lvert #1 \rvert}
\providecommand{\Pow}{\mathcal{P}}
\providecommand{\Setabs}[2]{\{\, #1 \mid #2 \,\}}
\providecommand{\comp}{\circ}
\providecommand{\Char}{\chi}
\providecommand{\Prf}{\operatorname{Prf}}
\providecommand{\Prov}{\operatorname{Prov}}
\providecommand{\Proj}{\operatorname{proj}}
\providecommand{\Gn}{\operatorname{gn}}
\providecommand{\gn}{\operatorname{gn}}
\providecommand{\Value}[1]{\operatorname{val}(#1)}
\providecommand{\Deduce}{\mathrm{D}}
\providecommand{\Dual}{\mathrm{d}}
\providecommand{\Ax}{\mathrm{Ax}}
\providecommand{\pSat}[3]{#1 \models #2\,[#3]}
\providecommand{\pAssign}{\sigma}
\providecommand{\bforall}{\forall}
\providecommand{\bexists}{\exists}
\providecommand{\liff}{\mathbin{\leftrightarrow}}
\providecommand{\FV}{\operatorname{FV}}
\providecommand{\UnaryInf}{\ensuremath{\mid}}
\providecommand{\BinaryInf}{\ensuremath{\mid}}
\providecommand{\depth}{\operatorname{depth}}
\providecommand{\oliflabeldef}[3]{#1}
\providecommand{\TMblank}{\sqcup}
\providecommand{\TMstroke}{\mid}
\providecommand{\TMstay}{\mathrm{S}}
\providecommand{\TMright}{\mathrm{R}}
\providecommand{\TMleft}{\mathrm{L}}
\providecommand{\TMendtape}{\triangleright}
\providecommand{\umin}{\mu}
\providecommand{\Sum}{\sum}
\providecommand{\citep}[1]{\cite{#1}}
\providecommand{\citet}[1]{\cite{#1}}
\providecommand{\citeauthor}[1]{\cite{#1}}
\providecommand{\citeyear}[1]{\cite{#1}}

% ===================== OLP compatibility layer (3) =====================
% --- third compatibility block: remaining notation seen in the frozen OLP units ---
\providecommand{\iso}{\cong}
\providecommand{\nszero}{\mathrm{0}}
\providecommand{\nssucc}{\mathrm{S}}
\providecommand{\nsplus}{\mathbin{+}}
\providecommand{\nssless}{<}
\providecommand{\nsless}{<}
\providecommand{\PL}{\mathrm{PL}}
\providecommand{\Nec}{\mathord{\Box}}
\providecommand{\RK}{\mathrm{RK}}
\providecommand{\ST}{\mathrm{ST}}
\providecommand{\SSubst}[3]{#1[#2/#3]}
\providecommand{\subst}[3]{#1[#2/#3]}
\providecommand{\tagfalse}{\mathrm{F}}
\providecommand{\cutr}{\mathrm{Cut}}
\providecommand{\cardneq}{\neq}
\providecommand{\issucc}[1]{\operatorname{succ}(#1)}
\providecommand{\tup}[1]{\langle #1 \rangle}
\providecommand{\Var}{\mathrm{Var}}
\providecommand{\Bew}{\operatorname{Bew}}
\providecommand{\Con}{\operatorname{Con}}
\providecommand{\Cons}{\operatorname{Cons}}
\providecommand{\pDeduce}{\mathrm{D}}
\providecommand{\psat}{\models}
\providecommand{\den}[1]{\llbracket #1 \rrbracket}
\providecommand{\Sem}[2]{\llbracket #1 \rrbracket_{#2}}
\providecommand{\DomainOf}[1]{\lvert #1 \rvert}
\providecommand{\Form}{\mathrm{Form}}
\providecommand{\Terms}{\mathrm{Term}}
\providecommand{\Atom}{\mathrm{Atom}}
\providecommand{\FVars}[1]{\operatorname{FV}(#1)}
\providecommand{\fv}[1]{\operatorname{fv}(#1)}
\providecommand{\tuplep}[1]{\langle #1 \rangle}
\providecommand{\restrict}[2]{#1 \upharpoonright #2}
\providecommand{\sto}{\rightarrow}
\providecommand{\FmlaSet}[1]{\mathrm{Fmla}_{#1}}
\providecommand{\Thm}{\mathrm{Thm}}
\providecommand{\Eval}{\operatorname{val}}
\providecommand{\blank}{\sqcup}
\providecommand{\blankTape}{\sqcup}
\providecommand{\Consis}{\operatorname{Con}}
\providecommand{\drs}{\mathrm{d}}
\providecommand{\VRule}{\mid}
\providecommand{\Bet}{\mathrm{B}}
\providecommand{\Seg}[1]{\mathrm{Seg}_{#1}}
\providecommand{\Enums}{\mathrm{En}}
\providecommand{\Z}{\mathbb{Z}}
\providecommand{\N}{\mathbb{N}}
\providecommand{\abs}[1]{\lvert #1 \rvert}
\providecommand{\floor}[1]{\lfloor #1 \rfloor}
\providecommand{\ceil}[1]{\lceil #1 \rceil}

% ===================== OLP compatibility layer (4) =====================
% --- fourth compatibility block: harvested from the per-candidate compile logs ---
\providecommand{\beredpar}{\mathrel{\rightarrow_\beta}}
\makeatletter
\def\bered{\@ifnextchar[{\bered@i}{\mathrel{\rightarrow_\beta}}}
\def\bered@i[#1]{\mathrel{\rightarrow_\beta}}
\def\bred{\@ifnextchar[{\bred@i}{\mathrel{\rightarrow_\beta}}}
\def\bred@i[#1]{\mathrel{\rightarrow_\beta}}
\makeatother
\providecommand{\beredone}{\mathrel{=_\beta}}
\providecommand{\bredone}{\mathrel{=_\beta}}
\providecommand{\xredone}{\mathrel{=_\beta}}
\providecommand{\aconvone}{\mathrel{=_\alpha}}
\providecommand{\cardless}{<}
\providecommand{\tagrefs}[1]{#1}
\providecommand{\ordplus}{\mathbin{+}}
\providecommand{\funrestrictionto}[2]{#1 \upharpoonright #2}
\providecommand{\fundefined}{\mathord{\uparrow}}
\providecommand{\Trm}{\mathrm{Trm}}
\providecommand{\OProv}{\operatorname{Prov}}
\providecommand{\OPrf}{\operatorname{Prf}}
\providecommand{\aeq}{\mathrel{=}}
\providecommand{\mClass}[1]{[#1]}
\providecommand{\equivclass}[1]{[#1]}
\providecommand{\pheight}{\operatorname{ht}}
\providecommand{\phy}{\operatorname{ht}}
\providecommand{\Cut}{\mathrm{Cut}}
\providecommand{\QuantRank}{\operatorname{qr}}
\providecommand{\elemequiv}{\mathrel{\equiv}}
\providecommand{\PIso}{\mathrm{PIso}}
\providecommand{\Theory}{\mathrm{Th}}
\providecommand{\supstrict}{\subsetneq}
\providecommand{\trcl}{\operatorname{trcl}}
\providecommand{\cref}[1]{\ref{#1}}
\providecommand{\Discharge}{\mathrm{Dis}}
\providecommand{\Weakening}{\mathrm{Wk}}
\providecommand{\ZFminus}{\mathrm{ZF}^{-}}
\makeatletter
\renewcommand{\Dual}{\@ifnextchar[{\Dual@i}{\ifmmode\mathrm{d}\else$\mathrm{d}$\fi}}
\def\Dual@i[#1]{\ifmmode\mathrm{d}[#1]\else$\mathrm{d}[#1]$\fi}
\makeatother
\providecommand{\RightR}{\mathrm{R}}
\providecommand{\LeftR}{\mathrm{L}}
\providecommand{\Intro}{\mathrm{I}}
\providecommand{\Axiom}{\mathrm{Ax}}
\providecommand{\formula}[1]{#1}
\providecommand{\stared}{\ast}
\providecommand{\oset}[2]{\{\, #1 \mid #2 \,\}}
\providecommand{\setof}[2]{\{\, #1 \mid #2 \,\}}
\providecommand{\pow}[1]{\mathcal{P}(#1)}
\providecommand{\ordpair}[2]{\langle #1, #2 \rangle}
\providecommand{\ordtriple}[3]{\langle #1, #2, #3 \rangle}
\providecommand{\union}{\cup}
\providecommand{\inter}{\cap}
\providecommand{\seq}[1]{\langle #1 \rangle}
\providecommand{\SuchThat}{\mid}
\providecommand{\where}{\mid}
\providecommand{\TVar}{\mathrm{TVar}}
\providecommand{\Bool}{\mathrm{Bool}}
\providecommand{\Val}{\operatorname{Val}}
\providecommand{\TL}[1]{\mathrm{TL}}
\providecommand{\xstared}{\ast}

% ===================== OLP compatibility layer (5) =====================
% --- fifth compatibility block: rescues the quarantined OLP candidates ---
\providecommand{\ident}{\operatorname{id}}
\providecommand{\setrank}{\operatorname{rank}}
\providecommand{\cardfont}{\mathfrak{c}}
\providecommand{\Complement}{\overline}
\providecommand{\QR}{\mathrm{QR}}
\providecommand{\becd}{\ensuremath{\ulcorner}}
\providecommand{\disjointsum}{\mathbin{\oplus}}
\providecommand{\cardeq}[2]{#1 = #2}
\providecommand{\xred}{\mathrel{\rightarrow}}
% --- sixth compatibility block: names reported by real compile logs ---
\providecommand{\ORefut}{\operatorname{Refut}}
\providecommand{\ORProv}{\operatorname{Prov}}
\providecommand{\cutrank}{\operatorname{cr}}
\providecommand{\Refut}{\operatorname{Refut}}
\providecommand{\DischargeRule}{\mathrm{Dis}}
\providecommand{\Contraction}{\mathrm{C}}
\providecommand{\bottomAlignProof}{\relax}
\providecommand{\FalseCl}{\bot}
\providecommand{\red}{\mathrel{\rightarrow}}
\providecommand{\varAssign}{\sigma}
\providecommand{\defiff}{\mathrel{\leftrightarrow}}
\providecommand{\FalseInt}{\bot}
\providecommand{\fn}{\operatorname{fn}}
\providecommand{\dotted}{\cdot}

% --- tag-macro semantic fix (2026-10-13, final): \iftag CHOOSES a branch and the OLP proof body is
% always the SECOND group (\iftag{tag}{body}, {tag}{body}{c}, {tag}{body}{c}{d}); a definition that
% prints the tag literal drops the proof body. Bodies contain blank lines, hence \long.
\makeatletter
\providecommand{\iftag}{}
\long\def\iftag#1#2{\@ifnextchar\bgroup{\iftag@three{#2}}{\iftag@stop{#2}}}
\long\def\iftag@three#1#2{\iftag@stop{#1}\@ifnextchar\bgroup{\iftag@eat}{\relax}}
\long\def\iftag@eat#1{}
\long\def\iftag@stop#1{#1}
\long\def\oliflabeldef#1#2{\@ifnextchar\bgroup{\@gobble}{\iftag@stop{#2}}}
\providecommand{\tagrefs}[1]{}
\providecommand{\settag}[1]{}
\makeatother
\begin{document}

% ==== OLP file: content/incompleteness/arithmetization-syntax/arithmetization-syntax.tex ====
% Part: incompleteness
% Chapter: arithmetization-syntax



\olchapter{inc}{art}{Arithmetization of Syntax}

\begin{editorial}
  Note that arithmetization for signed tableaux is not yet available.
\end{editorial}

\olimport{introduction}

\olimport{coding-symbols}

\olimport{coding-terms}

\olimport{coding-formulas}

\olimport{substitution}

\iftag{prfSC}{%
  \olimport{proofs-in-lk}
}{}

\iftag{prfND}{%
  \olimport{proofs-in-nd}
}{}

\iftag{prfAX}{
  \olimport{proofs-in-ax}
}{}

\OLEndChapterHook


% ==== OLP file: content/incompleteness/arithmetization-syntax/coding-formulas.tex ====
% Part: incompleteness
% Chapter: arithmetization-syntax
% Section: coding-formulas



\olfileid{inc}{art}{frm}
\olsection{Coding \printtoken{P}{formula}}

Once we have defined the relation $\fn{Term}(x)$ primitive
recursively, we can use it to define the corresponding relation for
!!{formula}s, $\fn{Frm}(x)$ primitive recursively.

\begin{prop}
The relation $\fn{Atom}(x)$ which holds iff $x$ is the G\"odel number
of an atomic !!{formula}, is primitive recursive.
\end{prop}

\begin{proof}
The number $x$ is the G\"odel number of an atomic !!{formula} iff
one of the following holds:
\begin{enumerate}
\item There are $n$, $j < x$, and $z < x$ such that for each $i < n$,
  $\fn{Term}((z)_i)$ and $x = $
\[
\Gn{\Obj P^n_j(} \concat \fn{flatten}(z) \concat \Gn{)}.
\]
\item There are $z_1, z_2 < x$ such that $\fn{Term}(z_1)$,
  $\fn{Term}(z_2)$, and $x = $
\[
\Gn{{\eq}(} \concat z_1 \concat \Gn{,} \concat z_2 \concat{\Gn{)}}.
\]
  \tagitem{prvFalse}{$x = \Gn{\lfalse}$.}{}
  \tagitem{prvTrue}{$x = \Gn{\ltrue}$.}{}
\end{enumerate}
\end{proof}

\begin{prop}
\ollabel{prop:frm-primrec}
The relation $\fn{Frm}(x)$ which holds iff $x$ is the G\"odel number
of !!a{formula} is primitive recursive.
\end{prop}

\begin{proof}
A sequence of symbols~$s$ is !!a{formula} iff there is formation
sequence~$s_0$, \dots, $s_{k-1} = s$ of !!{formula} which records
how~$s$ was formed from atomic !!{formula}s according to the
formation rules.  The code for each $s_i$ (and indeed of the code of
the sequence $\tuple{s_0, \dots, s_{k-1}}$) is less than the code~$x$
of~$s$.
\end{proof}

\begin{prob}
Give a detailed proof of \olref[inc][art][frm]{prop:frm-primrec} along
the lines of the first proof of
\olref[inc][art][trm]{prop:term-primrec}.
\end{prob}

\begin{prop}
\ollabel{prop:freeocc-primrec}
The relation $\fn{FreeOcc}(x, z, i)$, which holds iff the $i$-th
symbol of the formula with G\"odel number~$x$ is a free occurrence of
the variable with G\"odel number~$z$, is primitive recursive.
\end{prop}

\begin{proof}
Exercise.
\end{proof}

\begin{prob}
Prove \olref[inc][art][frm]{prop:freeocc-primrec}.  You may make use
of the fact that any substring of !!a{formula} which is !!a{formula}
is a sub-!!{formula} of it.
\end{prob}

\begin{prop}
The property $\fn{Sent}(x)$ which holds iff $x$~is the G\"odel
number of !!a{sentence} is primitive recursive.
\end{prop}

\begin{proof}
A !!{sentence} is !!a{formula} without free occurrences of
!!{variable}s. So $\fn{Sent}(x)$ holds iff
\begin{multline*}
\bforall{i<\len{x}}{\bforall{z<x}{}}\\
(\bexists{j<z}{z=\Gn{\Obj v_j}} \lif \lnot\fn{FreeOcc}(x,z,i)).
\end{multline*}
\end{proof}



% ==== OLP file: content/incompleteness/arithmetization-syntax/coding-symbols.tex ====
% Part: incompleteness
% Chapter: arithmetization-syntax
% Section: coding-symbols



\olfileid{inc}{art}{cod}
\olsection{Coding Symbols}

The basic language~$\Lang L$ of first order logic makes use of the symbols
\[
\lfalse \quad \lnot \quad \lor \quad \land \quad \lif \quad \lforall
\quad \lexists \quad \eq \quad ( \quad ) \quad ,
\]
together with !!{enumerable} sets of variables and !!{constant}s, and
!!{enumerable} sets of !!{function}s and !!{predicate}s of arbitrary
arity.  We can assign \emph{codes} to each of these symbols in such a
way that every symbol is assigned a unique number as its code, and no
two different symbols are assigned the same number.  We know that this
is possible since the set of all symbols is !!{enumerable} and so
there is !!a{bijection} between it and the set of natural numbers.
But we want to make sure that we can recover the symbol (as well as
some information about it, e.g., the arity of !!a{function}) from its
code in a computable way.  There are many possible ways of doing this,
of course.  Here is one such way, which uses primitive recursive
functions.  (Recall that $\tuple{n_0, \dots, n_k}$ is the number
coding the sequence of numbers $n_0$, \dots, $n_k$.)

\begin{defn}
If $s$ is a symbol of~$\Lang L$, let the \emph{symbol code}~$\scode s$ be
defined as follows:
\begin{enumerate}
\item If $s$ is among the logical symbols, $\scode s$ is given by the
  following table:
\[
\begin{array}{cccccccccc}
  \lfalse & \lnot & \lor & \land & \lif & \lforall \\
\tuple{0, 0} & \tuple{0, 1} & \tuple{0, 2} & \tuple{0, 3} &
\tuple{0, 4} & \tuple{0, 5} \\
\lexists & \eq & ( & ) & ,\\
\tuple{0, 6} & \tuple{0, 7} &
\tuple{0, 8} & \tuple{0, 9} & \tuple{0, 10}
\end{array}
\]
\item If $s$ is the $i$-th variable $\Obj v_i$, then $\scode s = \tuple{1, i}$.
\item If $s$ is the $i$-th !!{constant}~$\Obj c_i$, then
  $\scode s = \tuple{2, i}$.
\item If $s$ is the $i$-th $n$-ary !!{function}~$\Obj f_i^n$, then
  $\scode s = \tuple{3, n, i}$.
\item If $s$ is the $i$-th $n$-ary !!{predicate}~$\Obj P_i^n$, then
  $\scode s = \tuple{4, n, i}$.
\end{enumerate}
\end{defn}

\begin{prop}
The following relations are primitive recursive:
\begin{enumerate}
\item $\fn{Fn}(x, n)$ iff $x$ is the code of $\Obj f^n_i$ for
  some~$i$, i.e., $x$ is the code of an $n$-ary !!{function}.
\item $\fn{Pred}(x, n)$ iff $x$ is the code of $\Obj P^n_i$ for
  some~$i$ or $x$ is the code of $\eq$ and $n = 2$, i.e., $x$ is the
  code of an $n$-ary !!{predicate}.
\end{enumerate}
\end{prop}

\begin{defn}
If $s_0, \dots, s_{n-1}$ is a sequence of symbols, its \emph{G\"odel
  number} is $\tuple{\scode{s_0}, \dots, \scode{s_{n-1}}}$.
\end{defn}

\begin{explain}
Note that \emph{codes} and \emph{G\"odel numbers} are different
things. For instance, the variable~$\Obj v_5$ has a code~$\scode{\Obj
  v_5} = \tuple{1, 5} = 2^2\cdot 3^6$. But the variable~$\Obj v_5$
considered as a term is also a sequence of symbols (of
length~$1$). The \emph{G\"odel number}~$\Gn{\Obj v_5}$ of the
\emph{term}~$\Obj v_5$ is $\tuple{\scode{\Obj v_5}} = 2^{\scode{\Obj
    v_5} + 1} = 2^{2^2\cdot 3^6 + 1}$.
\end{explain}

\begin{ex}
Recall that if $k_0$, \dots, $k_{n-1}$ is a sequence of numbers, then
the code of the sequence $\tuple{k_0, \dots, k_{n-1}}$ in the
power-of-primes coding is
\[
2^{k_0+1}\cdot3^{k_1+1}\cdot \dots \cdot p_{n-1}^{k_{n-1}+1},
\]
where $p_i$ is the $i$-th prime (starting with $p_0 = 2$). So for
instance, the formula $\eq[\Obj v_0][\Obj 0]$, or, more explicitly,
${\eq}(\Obj v_0,\Obj c_0)$, has the G\"odel number
\[
\tuple{\scode{\eq},\scode{(},\scode{\Obj v_0},\scode{,}, \scode{\Obj
    c_0},\scode{)}}.
\]
Here, $\scode{\eq}$ is $\tuple{0,7} = 2^{0+1}\cdot
3^{7+1}$, $\scode{\Obj v_0}$ is $\tuple{1,0} = 2^{1+1}\cdot3^{0+1}$,
etc. So $\Gn{=(\Obj v_0,\Obj c_0)}$ is
\begin{multline*}
2^{\scode{=} + 1}\cdot 3^{\scode{(}+1}\cdot 5^{\scode{\Obj v_0}+1}
\cdot 7^{\scode{,} + 1} \cdot 11^{\scode{\Obj c_0}+1} \cdot
13^{\scode{)}+1} = \\
2^{2^1\cdot 3^8 + 1}\cdot 3^{2^1\cdot 3^9+1}\cdot 5^{2^2\cdot 3^1+1}
\cdot 7^{2^1\cdot 3^{11} + 1} \cdot 11^{2^3\cdot3^1+1} \cdot
13^{2^1\cdot3^{10}+1} = \\
2^{13\,123}\cdot 3^{39\,367}\cdot 5^{13}\cdot 7^{354\,295}\cdot11^{25}\cdot13^{118\,099}.
\end{multline*}
\end{ex}


% ==== OLP file: content/incompleteness/arithmetization-syntax/coding-terms.tex ====
% Part: incompleteness
% Chapter: arithmetization-syntax
% Section: coding-terms



\olfileid{inc}{art}{trm}
\olsection{Coding Terms}

\begin{explain}
A term is simply a certain kind of sequence of symbols: it is built up
inductively from constants and variables according to the formation
rules for terms.  Since sequences of symbols can be coded as
numbers---using a coding scheme for the symbols plus a way to code
sequences of numbers---assigning G\"odel numbers to terms is not
difficult.  The challenge is rather to show that the property a number
has if it is the G\"odel number of a correctly formed term is
computable, or in fact primitive recursive.
\end{explain}

!!^{variable}s and !!{constant}s are the simplest terms, and testing
whether $x$ is the G\"odel number of such a term is easy:
$\fn{Var}(x)$ holds if $x$ is $\Gn{\Obj v_i}$ for some~$i$. In other words,
$x$~is a sequence of length~$1$ and its single element $(x)_0$ is the
code of some !!{variable}~$\Obj v_i$, i.e., $x$ is $\tuple{\tuple{1, i}}$ for
some~$i$. Similarly, $\fn{Const}(x)$ holds if $x$ is $\Gn{\Obj c_i}$ for
some~$i$. Both of these relations are primitive recursive, since if
such an $i$ exists, it must be $< x$:
\begin{align*}
  \fn{Var}(x) & \defiff \bexists{i<x}{x = \tuple{\tuple{1, i}}}\\
  \fn{Const}(x) & \defiff \bexists{i<x}{x = \tuple{\tuple{2, i}}}
\end{align*}

\begin{prop}
\ollabel{prop:term-primrec}
The relations $\fn{Term}(x)$ and $\fn{ClTerm}(x)$ which hold iff $x$
is the G\"odel number of a term or a closed term, respectively, are
primitive recursive.
\end{prop}

\begin{proof}
A sequence of symbols~$s$ is a term iff there is a sequence~$s_0$,
\dots, $s_{k-1} = s$ of terms which records how the term~$s$ was formed
from !!{constant}s and !!{variable}s according to the formation rules
for terms. To express that such a putative formation sequence follows
the formation rules it has to be the case that, for each $i < k$, either
\begin{enumerate}
\item $s_i$ is !!a{variable}~$\Obj v_j$, or
\item $s_i$ is !!a{constant}~$\Obj c_j$, or
\item $s_i$ is built from $n$ terms $t_1$, \dots, $t_n$ occurring
  prior to place~$i$ using an $n$-place !!{function}~$\Obj f^n_j$.
\end{enumerate}
To show that the corresponding relation on G\"odel numbers is
primitive recursive, we have to express this condition primitive
recursively, i.e., using primitive recursive functions, relations, and
bounded quantification.

Suppose $y$ is the number that codes the sequence $s_0$, \dots, $s_{k-1}$,
i.e., $y = \tuple{\Gn{s_0}, \dots, \Gn{s_{k-1}}}$.  It codes a formation
sequence for the term with G\"odel number~$x$ iff for all $i < k$:
\begin{enumerate}
\item $\fn{Var}((y)_i)$, or
\item $\fn{Const}((y)_i)$, or
\item there is an $n$ and a number~$z = \tuple{z_1, \dots, z_n}$ such
  that each $z_l$ is equal to some $(y)_{i'}$ for $i' < i$ and
\[
(y)_i = \Gn{\Obj f^n_j(} \concat \fn{flatten}(z) \concat \Gn{)},
\]
\end{enumerate}
and moreover $(y)_{k-1} = x$.  (The function $\fn{flatten}(z)$ turns
the sequence $\tuple{\Gn{t_1}, \dots, \Gn{t_n}}$ into $\Gn{t_1, \dots,
  t_n}$ and is primitive recursive.)

The indices $j$, $n$, the G\"odel numbers $z_l$ of the terms $t_l$,
and the code~$z$ of the sequence~$\tuple{z_1, \dots, z_n}$, in (3) are
all less than~$y$.  We can replace $k$ above with $\len{y}$.  Hence we
can express ``$y$ is the code of a formation sequence of the term with
G\"odel number~$x$'' in a way that shows that this relation is
primitive recursive.

We now just have to convince ourselves that there is a primitive
recursive bound on~$y$.  But if $x$ is the G\"odel number of a term,
it must have a formation sequence with at most $\len{x}$ terms (since
every term in the formation sequence of~$s$ must start at some place
in~$s$, and no two subterms can start at the same place).  The G\"odel
number of each subterm of~$s$ is of course $\le x$.  Hence, there
always is a formation sequence with code $\le p_{k-1}^{k(x+1)}$, where
$k=\len{x}$.

For $\fn{ClTerm}$, simply leave out the clause for !!{variable}s.
\end{proof}

\begin{prob}
Show that the function $\fn{flatten}(z)$, which turns the sequence
$\tuple{\Gn{t_1}, \dots, \Gn{t_n}}$ into $\Gn{t_1, \dots, t_n}$, is
primitive recursive.
\end{prob}

\begin{prop}\ollabel{prop:num-primrec}
  The function $\fn{num}(n) = \Gn{\num{n}}$ is primitive recursive.
\end{prop}

\begin{proof}
  We define $\fn{num}(n)$ by primitive recursion:
  \begin{align*}
    \fn{num}(0) & = \Gn{\Obj 0}\\
    \fn{num}(n+1) & = \Gn{\prime(} \concat \fn{num}(n) \concat \Gn{)}.
  \end{align*}
\end{proof}


% ==== OLP file: content/incompleteness/arithmetization-syntax/introduction.tex ====
% Part: incompleteness
% Chapter: arithmetization-syntax
% Section: introduction



\olfileid{inc}{art}{int}
\olsection{Introduction}

In order to connect computability and logic, we need a way to talk
about the objects of logic (symbols, terms, !!{formula}s,
!!{derivation}s), operations on them, and their properties and
relations, in a way amenable to computational treatment.  We can do
this directly, by considering computable functions and relations on
symbols, sequences of symbols, and other objects built from them.
Since the objects of logical syntax are all finite and built from
!!a{enumerable} sets of symbols, this is possible for some models of
computation.  But other models of computation---such as the recursive
functions----are restricted to numbers, their relations and functions.
Moreover, ultimately we also want to be able to deal with syntax
within certain theories, specifically, in theories formulated in the
language of arithmetic.  In these cases it is necessary to
\emph{arithmetize} syntax, i.e., to represent syntactic objects,
operations on them, and their relations, as numbers, arithmetical
functions, and arithmetical relations, respectively. The idea, which
goes back to Leibniz, is to assign numbers to syntactic objects.

It is relatively straightforward to assign numbers to symbols as their
``codes.''  Some symbols pose a bit of a challenge, since, e.g., there
are infinitely many !!{variable}s, and even infinitely many
!!{function}s of each arity~$n$. But of course it's possible to assign
numbers to symbols systematically in such a way that, say, $\Obj v_2$
and $\Obj v_3$ are assigned different codes. Sequences of symbols
(such as terms and !!{formula}s) are a bigger challenge. But if we can
deal with sequences of numbers purely arithmetically (e.g., by the
powers-of-primes coding of sequences), we can extend the coding of
individual symbols to coding of sequences of symbols, and then further
to sequences or other arrangements of !!{formula}s, such as
!!{derivation}s. This extended coding is called ``G\"odel numbering.''
Every term, !!{formula}, and !!{derivation} is assigned a G\"odel
number.

By coding sequences of symbols as sequences of their codes, and by
choosing a system of coding sequences that can be dealt with using
computable functions, we can then also deal with G\"odel numbers using
computable functions.  In practice, all the relevant functions will be
primitive recursive.  For instance, computing the length of a sequence
and computing the $i$-th element of a sequence from the code of the
sequence are both primitive recursive. If the number coding the
sequence is, e.g., the G\"odel number of !!a{formula}~$!A$, we
immediately see that the length of !!a{formula} and the (code of the)
$i$-th symbol in !!a{formula} can also be computed from the G\"odel
number of~$!A$. It is a bit harder to prove that, e.g., the property
of being the G\"odel number of a correctly formed term or of a correct
!!{derivation} is primitive recursive.  It is nevertheless possible,
because the sequences of interest (terms, !!{formula}s,
!!{derivation}s) are inductively defined.

As an example, consider the operation of substitution. If $!A$ is a
formula, $x$ a variable, and $t$ a term, then $\Subst{!A}{t}{x}$ is
the result of replacing every free occurrence of~$x$ in~$!A$ by~$t$.
Now suppose we have assigned G\"odel numbers to $!A$, $x$, $t$---say,
$k$, $l$, and $m$, respectively.  The same scheme assigns a G\"odel
number to $\Subst{!A}{t}{x}$, say,~$n$.  This mapping---of $k$, $l$,
and $m$ to $n$---is the arithmetical analog of the substitution
operation. When the substitution operation maps $!A$, $x$, $t$ to
$\Subst{!A}{t}{x}$, the arithmetized substitution functions maps the
G\"odel numbers $k$, $l$, $m$ to the G\"odel number~$n$.  We will see
that this function is primitive recursive.

Arithmetization of syntax is not just of abstract interest, although
it was originally a non-trivial insight that languages like the
language of arithmetic, which do not come with mechanisms for
``talking about'' languages can, after all, formalize complex
properties of expressions.  It is then just a small step to ask what a
theory in this language, such as Peano arithmetic, can \emph{prove}
about its own language (including, e.g., whether !!{sentence}s are
provable or true).  This leads us to the famous limitative theorems of
G\"odel (about unprovability) and Tarski (the undefinability of
truth). But the trick of arithmetizing syntax is also important in
order to prove some important results in computability theory, e.g.,
about the computational power of theories or the relationship between
different models of computability.  The arithmetization of syntax
serves as a model for arithmetizing other objects and properties. For
instance, it is similarly possible to arithmetize configurations and
computations (say, of Turing machines). This makes it possible to
simulate computations in one model (e.g., Turing machines) in another
(e.g., recursive functions).


% ==== OLP file: content/incompleteness/arithmetization-syntax/proofs-in-ax.tex ====
% Part: incompleteness
% Chapter: arithmetization-syntax
% Section: proofs-in-ax



\olfileid{inc}{art}{pax}
\olsection{Axiomatic \usetoken{P}{derivation}}

\begin{explain}
In order to arithmetize axiomatic !!{derivation}s, we must represent
!!{derivation}s as numbers. Since !!{derivation}s are simply sequences
of !!{formula}s, the obvious approach is to code every !!{derivation}
as the code of the sequence of codes of !!{formula}s in it.
\end{explain}

\begin{defn}
If $\delta$ is an axiomatic !!{derivation} consisting of !!{formula}s
$!A_1$, \dots,~$!A_n$, then $\Gn{\delta}$ is
\[
\tuple{\Gn{!A_1}, \dots, \Gn{!A_n}}.
\]
\end{defn}

\begin{ex}
  Consider the very simple !!{derivation}:
  \begin{derivation}
  1. & $!B \lif (!B \lor !A)$ \\
  2. & $(!B \lif (!B \lor !A)) \lif (!A  \lif (!B \lif (!B \lor !A)))$\\
  3. & $!A  \lif (!B \lif (!B \lor !A))$
  \end{derivation}
  The G\"odel number of this !!{derivation} would be
  \begin{align*}
  \openTuple\, 
    & \Gn{!B \lif (!B \lor !A)}, \\
    &\Gn{(!B \lif (!B \lor !A)) \lif (!A  \lif (!B \lif (!B \lor !A)))},\\
    & \Gn{!A  \lif (!B \lif (!B \lor !A))} \,\closeTuple.
  \end{align*}
\end{ex}

\begin{explain}
Having settled on a representation of !!{derivation}s, we must also
show that we can manipulate such !!{derivation}s primitive recursively, and
express their essential properties and relations so.  Some operations
are simple: e.g., given a G\"odel number~$d$ of !!a{derivation},
$(d)_{\len{d}-1}$ gives us the G\"odel number of its end-!!{formula}.  Some are
much harder.  We'll at least sketch how to do this.  The goal is to
show that the relation ``$\delta$ is !!a{derivation} of~$!A$
from~$\Gamma$'' is primitive recursive in the G\"odel numbers of
$\delta$ and~$!A$.
\end{explain}

\begin{prop}
\ollabel{prop:followsby}
The following relations are primitive recursive:
\begin{enumerate}
\item $!A$ is an axiom.
\item The $i$-th line in $\delta$ is justified by modus ponens
\item The $i$-th line in $\delta$ is justified by \QR.
\item $\delta$ is a correct !!{derivation}.
\end{enumerate}
\end{prop}

\begin{proof}
We have to show that the corresponding relations between G\"odel
numbers of !!{formula}s and  G\"odel numbers of !!{derivation}s 
are primitive recursive.
\begin{enumerate}
\item We have a given list of axiom schemas, and $!A$ is an axiom if
  it is of the form given by one of these schemas.  Since the list of
  schemas is finite, it suffices to show that we can test primitive
  recursively, for each axiom schema, if $!A$ is of that form. For
  instance, consider the axiom schema
  \[
  !B \lif (!C \lif !B).
  \]
  $!A$ is an instance of this axiom schema if there are !!{formula}s
  $!B$ and $!C$ such that we obtain $!A$ when we concatenate `$($'
  with $!B$ with `$\lif$' with `$($' with $!C$ with `$\lif$' with $!B$
  and with~`$))$'. We can test the corresponding property of the
  G\"odel number $n$ of~$!A$, since concatenation of sequences is
  primitive recursive and the G\"odel numbers of $!B$ and $!C$ must
  be smaller than the G\"odel number of~$!A$, since when the relation
  holds, both $!B$ and $!C$ are sub-!!{formula}s of~$!A$. Hence, we
  can define:
  \begin{multline*}
  \fn{IsAx}_{!B \lif (!C \lif !B)}(n) \defiff \bexists{b< n}{
    \bexists{c < n}{(\fn{Sent}(b) \land \fn{Sent}(c) \land {}}}\\ n =
  \Gn{(} \concat b \concat \Gn{\lif} \concat \Gn{(} \concat c \concat
  \Gn{\lif} \concat b \concat \Gn{))}).
  \end{multline*}
  If we have such a definition for each axiom schema, their
  disjunction defines the property $\fn{IsAx}(n)$, ``$n$~is the
  G\"odel number of an axiom.''
\item The $i$-th line in $\delta$ is justified by modus ponens iff there
  are lines $j$ and $k < i$ where the !!{sentence} on line~$j$ is some
  formula $!A$, the sentence on line~$k$ is $!A \lif !B$, and the
  sentence on line~$i$ is~$!B$.
  \begin{multline*}
    \fn{MP}(d, i) \defiff \bexists{j < i}{\bexists{k < i}{}}\\
    (d)_k = \Gn{(} \concat (d)_j
      \concat \Gn{\lif} \concat (d)_i \concat \Gn{)}
  \end{multline*}
  Since bounded quantification, concatenation, and $=$ are primitive
  recursive, this defines a primitive recursive relation. 
\item A line in $\delta$ is justified by \QR{} if it is of the form
  $!B \lif \lforall[x][!A(x)]$, a preceding line is $!B \lif !A(c)$
  for some !!{constant}~$c$, and $c$ does on occur in~$!B$. This is
  the case iff
  \begin{enumerate}
  \item there is !!a{sentence}~$!B$ and
  \item !!a{formula}~$!A(x)$ with a single variable~$x$ free so that
  \item line $i$ contains $!B \lif \lforall[x][!A(x)]$
  \item some line $j < i$ contains $!B \lif \Subst{!A}{c}{x}$ for a constant~$c$
  \item which does not occur in~$!B$.
  \end{enumerate}
  All of these can be tested primitive recursively, since the G\"odel
  numbers of $!B$, $!A(x)$, and $x$ are less than the G\"odel number of the
  formula on line~$i$, and that of $a$ less than the G\"odel number of the
  formula on line~$j$:
  \begin{multline*}
    \fn{QR}_1(d, i) \defiff \bexists{b < (d)_i}{\bexists{x <
        (d)_i}{\bexists{a < (d)_i}{\bexists{c < (d)_j}{(}}}}
    \\ \fn{Var}(x) \land \fn{Const}(c) \land {} \\
    (d)_i = \Gn{(} \concat
    b \concat \Gn{\lif} \concat \Gn{\lforall} \concat x \concat a
    \concat \Gn{)} \land {}\\
    (d)_j = \Gn{(} \concat b \concat
    \Gn{\lif} \concat \fn{Subst}(a,c,x) \concat \Gn{)} \land {}\\ \fn{Sent}(b)
    \land \fn{Sent}(\fn{Subst}(a,c,x)) \land {} \bforall{k <
      \len{b}}{(b)_k \neq (c)_0})
  \end{multline*}
  Here we assume that $c$ and $x$ are the G\"odel numbers of the
  variable and constant considered as terms (i.e., not their symbol
  codes).  We test that $x$ is the only free variable of~$!A(x)$ by
  testing if $\Subst{!A(x)}{c}{x}$ is !!a{sentence}, and ensure that
  $c$ does not occur in $!B$ by requiring that every symbol of~$!B$ is
  different from~$c$.

  We leave the other version of \QR{} as an exercise.
\item $d$ is the G\"odel number of a correct !!{derivation} iff every
  line in it is an axiom, or justified by modus ponens or~\QR. Hence:
  \[
  \fn{Deriv}(d) \defiff \bforall{i < \len{d}}{(\fn{IsAx}((d)_i) \lor
    \fn{MP}(d,i) \lor \fn{QR}(d, i))}
  \]
\end{enumerate}
\end{proof}

\begin{prob}
Define the following relations as in
\olref[inc][art][pax]{prop:followsby}:
\begin{enumerate}
\item $\fn{IsAx}_{!A \lif (!B \lif (!A \land !B))}(n)$,
\item $\fn{IsAx}_{\lforall[x][!A(x)] \lif !A(t)}(n)$,
\item $\fn{QR}_{2}(d, i)$ (for the other version of \QR).
\end{enumerate}
\end{prob}

\begin{prop}
Suppose $\Gamma$ is a primitive recursive set of !!{sentence}s.  Then
the relation $\Prf[\Gamma](x, y)$ expressing ``$x$ is the code of
!!a{derivation}~$\delta$ of $!A$ from~$\Gamma$ and $y$ is the G\"odel
number of~$!A$'' is primitive recursive.
\end{prop}

\begin{proof}
Suppose ``$y \in \Gamma$'' is given by the primitive recursive
predicate~$R_\Gamma(y)$.  We have to show that the relation
$\Prf[\Gamma](x, y)$ is primitive recursive, where $\Prf[\Gamma](x,
y)$ holds iff $y$ is the G\"odel number of !!a{sentence}~$!A$ and
$x$~is the code of !!a{derivation} of~$!A$ from $\Gamma$.

By the previous proposition, the property $\fn{Deriv}(x)$ which holds
iff $x$ is the code of a correct !!{derivation}~$\delta$ is primitive
recursive. However, that definition did not take into account the set
$\Gamma$ as an additional way to justify lines in the !!{derivation}. Our
primitive recursive test of whether a line is justified by \QR{} also
left out of consideration the requirement that the constant~$c$ is not
allowed to occur in~$\Gamma$. It is possible to amend our definition so
that it takes into account $\Gamma$ directly, but it is easier to use
$\fn{Deriv}$ and the deduction theorem. $\Gamma \Proves !A$ iff
there is some finite list of !!{sentence}s $!B_1$, \dots, $!B_n \in
\Gamma$ such that $\{!B_1, \dots, !B_n\} \Proves !A$. And by the
deduction theorem, this is the case if $\Proves (!B_1 \lif (!B_2 \lif
\cdots (!B_n \lif !A)\cdots))$. Whether !!a{sentence} with G\"odel
number~$z$ is of this form can be tested primitive recursively.  So,
instead of considering $x$ as the G\"odel number of !!a{derivation} of
the !!{sentence} with G\"odel number~$y$ \emph{from $\Gamma$}, we consider
$x$ as the G\"odel number of !!a{derivation} of a nested conditional
of the above form from~$\emptyset$.

First, if we have a sequence of !!{sentence}s, we can primitive
recursively form the conditional with all these sentences as
antecedents and given !!{sentence} as consequent:
\begin{align*}
  \fn{hCond}(s, y, 0) & = y \\
  \fn{hCond}(s, y, n+1) & = \Gn{(} \concat (s)_{n} \concat \Gn{\lif}
  \concat \fn{Cond}(s, y, n) \concat \Gn{)}\\
  \fn{Cond}(s, y) & = \fn{hCond}(s, y, \len{s})\\
  \intertext{So we can define $\Prf[\Gamma](x, y)$ by}
  \Prf[\Gamma](x, y) & \defiff \bexists{s < \fn{sequenceBound}(x,x)}{(} \\
  &\qquad (x)_{\len{x}-1} = \fn{Cond}(s,y) \land {} \\
  &\qquad\bforall{i<\len{s}}{(s)_i \in \Gamma} \land {} \\
  &\qquad\fn{Deriv}(x)).
\end{align*}
The bound on~$s$ is given by considering that each $(s)_i$ is the
G\"odel number of a sub-!!{formula} of the last line of the !!{derivation},
i.e., is less than $(x)_{\len{x}-1}$. The number of antecedents $!B
\in \Gamma$, i.e., the length of~$s$, is less than the length of the
last line of~$x$.
\end{proof}


% ==== OLP file: content/incompleteness/arithmetization-syntax/proofs-in-lk.tex ====
% Part: incompleteness
% Chapter: arithmetization-syntax
% Section: proofs-in-lk



\olfileid{inc}{art}{plk}
\olsection{\usetoken{P}{derivation} in $\Log{LK}$}

\begin{explain}
In order to arithmetize !!{derivation}s, we must represent
!!{derivation}s as numbers. Since !!{derivation}s are trees of sequents
where each inference carries also a label, a recursive
representation is the most obvious approach: we represent a
!!{derivation} as a tuple, the components of which are the
end-sequent, the label, and the representations of the
sub-!!{derivation}s leading to the premises of the last inference.
\end{explain}

\begin{defn}
If $\Gamma$ is a finite sequence of !!{sentence}s, $\Gamma =
\tuple{!A_1, \dots, !A_n}$, then $\Gn{\Gamma} = \tuple{\Gn{!A_1},
  \dots, \Gn{!A_n}}$.

If $\Gamma \Sequent \Delta$ is a sequent, then a G\"odel number of
$\Gamma \Sequent \Delta$ is
\[
\Gn{\Gamma \Sequent \Delta} = \tuple{\Gn{\Gamma}, \Gn{\Delta}}
\]

If $\pi$ is !!a{derivation} in $\Log{LK}$, then $\Gn{\pi}$ is defined
as follows:
\begin{enumerate}
\item If $\pi$ consists only
  of the initial sequent $\Gamma \Sequent \Delta$, then $\Gn{\pi}$ is 
  \[
    \tuple{0, \Gn{\Gamma \Sequent \Delta}}.
  \]
\item If $\pi$ ends in an inference with one or two premises, has
  $\Gamma \Sequent \Delta$ as its conclusion, and $\pi_1$ and $\pi_2$ are the immediate subproof ending
  in the premise of the last inference, then $\Gn{\pi}$ is 
  \begin{align*}
    & \tuple{1, \Gn{\pi_1}, \Gn{\Gamma \Sequent \Delta}, k} \text{ or}\\ 
    & \tuple{2, \Gn{\pi_1}, \Gn{\pi_2}, \Gn{\Gamma \Sequent \Delta}, k}, 
  \end{align*}
  respectively, where $k$ is given by the following table according to
  which rule was used in the last inference:

  \begin{tabular}{lccccccc}
    \text{Rule:} & \LeftR{\Weakening} & \RightR{\Weakening} &
    \LeftR{\Contraction} & \RightR{\Contraction} &
    \LeftR{\Exchange} & \RightR{\Exchange} \\
    $k$: & 1 & 2 & 3 & 4 & 5 & 6 \\[2ex]
    \text{Rule:} & \LeftR{\lnot} & \RightR{\lnot} &
    \LeftR{\land} & \RightR{\land} & 
    \LeftR{\lor} & \RightR{\lor} \\
    $k$: & 7 & 8 & 9 & 10 & 11 & 12 \\[2ex]
    \text{Rule:} & \LeftR{\lif} & \RightR{\lif} &
    \LeftR{\lforall} & \RightR{\lforall} &
    \LeftR{\lexists} & \RightR{\lexists} \\
    $k$: & 13 & 14 & 15 & 16 & 17 & 18 \\[2ex]
    \text{Rule:} & \Cut & = \\
    $k$: & 19 & 20
  \end{tabular}
\end{enumerate}
\end{defn}

\begin{ex}
  Consider the very simple !!{derivation}
  \begin{prooftree}
    \Axiom$!A \fCenter !A$
    \RightLabel{\LeftR{\land}}
    \UnaryInf$!A \land !B \fCenter !A$
    \RightLabel{\RightR{\lif}}
    \UnaryInf$\fCenter (!A \land !B) \lif !A$
  \end{prooftree}
  The G\"odel number of the initial sequent would be $p_0 = \tuple{0,
    \Gn{!A \Sequent !A}}$.  The G\"odel number of the !!{derivation}
    ending in the conclusion of~$\LeftR{\land}$ would be $p_1 =
    \tuple{1, p_0, \Gn{!A \land !B \Sequent !A}, 9}$ ($1$ since
    $\LeftR{\land}$ has one premise, the G\"odel number of the
    conclusion~$!A \land !B \Sequent !A$, and $9$ is the number coding
    $\LeftR{\land}$). The G\"odel number of the entire !!{derivation} then
    is $\tuple{1, p_1, \Gn{\Sequent (!A \land !B) \lif !A)}, 14}$, i.e.,
  \[
  \tuple{1, \tuple{1, \tuple{0, \Gn{!A \Sequent !A)}}, \Gn{!A \land !B \Sequent !A}, 9},
    \Gn{\Sequent (!A \land !B) \lif !A}, 14}.
  \]
\end{ex}

\begin{explain}
Having settled on a representation of !!{derivation}s, we must also
show that we can manipulate such !!{derivation}s primitive recursively,
and express their essential properties and relations so.  Some
operations are simple: e.g., given a G\"odel number~$p$ of
!!a{derivation}, $\fn{EndSeq}(p) = (p)_{(p)_0+1}$ gives us the G\"odel
number of its end-sequent and $\fn{LastRule}(p) = (p)_{(p)_0+2}$ the
code of its last rule.  The property $\fn{Sequent}(s)$ defined by
\[
  \len{s} = 2 \land \bforall{i<\len{(s)_0} + \len{(s)_1}}{\fn{Sent}(((s)_0 \concat (s)_1)_i)}
\]
holds of $s$ iff $s$ is the G\"odel number of a sequent consisting of
!!{sentence}s. Some are much harder.  We'll at least sketch how to do
this.  The goal is to show that the relation ``$\pi$ is
!!a{derivation} of~$!A$ from~$\Gamma$'' is a primitive recursive
relation of the G\"odel numbers of $\pi$ and~$!A$.
\end{explain}

\begin{prop}\ollabel{prop:followsby}
  The property $\fn{Correct}(p)$ which holds iff the last inference in
  the !!{derivation}~$\pi$ with G\"odel number~$p$ is correct, is
  primitive recursive.
\end{prop}

\begin{proof}
  $\Gamma \Sequent \Delta$ is an initial sequent if either there
  is !!a{sentence}~$!A$ such that $\Gamma \Sequent \Delta$ is $!A
  \Sequent !A$, or there is a term~$t$ such that $\Gamma \Sequent
  \Delta$ is $\emptyset \Sequent \eq[t][t]$.  In terms of G\"odel
  numbers, $\fn{InitSeq}(s)$ holds iff
  \begin{align*}
  \bexists{x < s}{} (\fn{Sent}(x) & \land
  s = \tuple{\tuple{x},\tuple{x}})
  \lor {}\\
  \bexists{t<s}{} (\fn{Term}(t) & \land
  s = \tuple{0, \tuple{\Gn{{\eq}(} \concat t \concat \Gn{,} \concat t \concat \Gn{)}}}).
  \end{align*}

  We also have to show that for each rule of inference~$R$ the
  relation $\fn{FollowsBy}_R(p)$ is primitive recursive, where
  $\fn{FollowsBy}_R(p)$ holds iff $p$ is the G\"odel number of
  !!{derivation}~$\pi$, and the end-sequent of~$\pi$ follows
  by a correct application of~$R$ from the immediate
  sub-!!{derivation}s of~$\pi$.

  A simple case is that of the \RightR{\land} rule. If $\pi$ ends in
  a correct $\RightR{\land}$ inference, it looks like this:
  \begin{prooftree}
    \AxiomC{}
    \RightLabel{$\pi_1$}
    \Deduce$\Gamma \fCenter \Delta, !A$
    
    \AxiomC{}
    \RightLabel{$\pi_2$}
    \Deduce$\Gamma \fCenter \Delta, !B$

    \RightLabel{\RightR\land}
    \BinaryInf$\Gamma \fCenter \Delta, !A \land !B$
  \end{prooftree}
  So, the last inference in the !!{derivation} $\pi$
  is a correct application of $\RightR{\land}$ iff
  there are sequences of !!{sentence}s $\Gamma$ and $\Delta$
  as well as two !!{sentence}s $!A$ and~$!B$ such that
  the end-sequent of $\pi_1$ is $\Gamma \Sequent \Delta, !A$,
  the end-sequent of $\pi_2$ is $\Gamma \Sequent \Delta, !B$,
  and the end-sequent of $\pi$ is $\Gamma \Sequent \Delta, !A \land !B$.
  We just have to translate this into G\"odel
  numbers.  If $s = \Gn{\Gamma \Sequent \Delta}$ then $(s)_0 =
  \Gn{\Gamma}$ and $(s)_1 = \Gn{\Delta}$.  So,
  $\fn{FollowsBy}_{\RightR{\land}}(p)$ holds iff
\begin{align*}
& \bexists{g < p}{\bexists{d < p}{\bexists{a < p}{\bexists{b < p}{\quad}}}} \\
& \qquad \fn{EndSequent}(p) = 
  \tuple{g, d \concat 
    \tuple{\Gn{(} \concat a \concat \Gn{\land} \concat b \concat \Gn{)}}} \land {} \\
& \qquad \fn{EndSequent}((p)_1) = \tuple{g, d \concat \tuple{a}} \land {}\\
& \qquad \fn{EndSequent}((p)_2) = \tuple{g, d \concat \tuple{b}} \land {}\\
& \qquad (p)_0 = 2 \land \fn{LastRule}(p) = 10.
\end{align*}
The individual lines express, respectively, ``there is a
sequence~($\Gamma$) with G\"odel number~$g$, there is a
sequence~($\Delta$) with G\"odel number~$d$, !!a{formula}~($!A$) with
G\"odel number~$a$, and !!a{formula}~($!B$) with G\"odel number~$b$,''
such that ``the end-sequent of $\pi$ is $\Gamma \Sequent \Delta, !A
\land !B$,'' ``the end-sequent of $\pi_1$ is $\Gamma \Sequent \Delta,
!A$,'' ``the end-sequent of $\pi_2$ is $\Gamma \Sequent \Delta, !B$,''
and ``$\pi$ has two immediate subderivations and the last inference
rule is $\RightR\land$ (with number~$10$).''

The last inference in~$\pi$ is a correct application of
$\RightR\lexists$ iff there are sequences $\Gamma$ and $\Delta$,
!!a{formula}~$!A$, a variable~$x$, and a term~$t$, such that
the end-sequent of $\pi$ is $\Gamma \Sequent \Delta, \lexists[x][!A]$
and the end-sequent of $\pi_1$ is $\Gamma \Sequent \Delta,
\Subst{!A}{t}{x}$. So in terms of G\"odel numbers, we have
$\fn{FollowsBy}_{\RightR\lexists}(p)$
iff
\begin{align*}
  & \bexists{g<p}{\bexists{d<p}{\bexists{a<p}{\bexists{x<p}{\bexists{t<p}{\quad}}}}}\\
  & \qquad \fn{EndSequent}(p) = 
  \tuple{
    g,
    d \concat \tuple{\Gn{\lexists} \concat x \concat a}
  } \land {}\\
  & \qquad \fn{EndSequent}((p)_1) = 
  \tuple{
    g,
    d \concat \tuple{\fn{Subst}(a, t, x)}
  } \land {}\\
  & \qquad (p)_0 = 1 \land \fn{LastRule}(p) = 18.
\end{align*}
We then define $\fn{Correct}(p)$ as
\begin{multline*}
  \fn{Sequent}(\fn{EndSequent}(p)) \land {}\\ 
  [(\fn{LastRule}(p) = 1 \land 
  \fn{FollowsBy}_{\LeftR\Weakening}(p)) \lor \dots \lor {}\\
  (\fn{LastRule}(p) = 20 \land \fn{FollowsBy}_{\eq}(p)) \lor {}\\
  (p)_0 = 0 \land \fn{InitialSeq}(\fn{EndSequent}(p))]
\end{multline*}
The first line ensures that the end-sequent of~$d$ is actually a
sequent consisting of !!{sentence}s. The last line covers the case
where $p$ is just an initial sequent.
\end{proof}

\begin{prob}
  Define the following properties as in
  \olref[inc][art][plk]{prop:followsby}:
  \begin{enumerate}
  \item $\fn{FollowsBy}_{\Cut}(p)$,
  \item $\fn{FollowsBy}_{\LeftR\lif}(p)$,
  \item $\fn{FollowsBy}_{\eq}(p)$,
  \item $\fn{FollowsBy}_{\RightR{\lforall}}(p)$.
  \end{enumerate}
  For the last one, you will have to also show that you can test
  primitive recursively if the last inference of the
  !!{derivation} with G\"odel number~$p$ satisfies the eigenvariable
  condition, i.e., the eigenvariable~$a$ of the $\RightR{\lforall}$
  does not occur in the end-sequent.
  \end{prob}
  

\begin{prop}
  \ollabel{prop:deriv}
  The relation $\fn{Deriv}(p)$ which holds if $p$ is the G\"odel
  number of a correct !!{derivation}~$\pi$, is primitive recursive.
\end{prop}

\begin{proof}
  !!^a{derivation}~$\pi$ is correct if every one of its inferences
  is a correct application of a rule, i.e., if every one of its
  sub-!!{derivation}s ends in a correct inference. So, $\fn{Deriv}(d)$
  iff
  \[
  \bforall{i<\len{\fn{SubtreeSeq}(p)}}{\fn{Correct}((\fn{SubtreeSeq}(p))_i}.
  \]
\end{proof}


\begin{prop}
Suppose $\Gamma$ is a primitive recursive set of !!{sentence}s.  Then
the relation $\Prf[\Gamma](x, y)$ expressing ``$x$ is the code of
!!a{derivation}~$\pi$ of $\Gamma_0 \Sequent !A$ for some finite
$\Gamma_0 \subseteq \Gamma$ and $y$ is the G\"odel number of~$!A$'' is
primitive recursive.
\end{prop}

\begin{proof}
Suppose ``$y \in \Gamma$'' is given by the primitive recursive
predicate~$R_\Gamma(y)$.  We have to show that $\Prf[\Gamma](x, y)$
which holds iff $y$ is the G\"odel number of a sentence~$!A$ and
$x$~is the code of an $\Log{LK}$-!!{derivation} with end-sequent
$\Gamma_0 \Sequent !A$ is primitive recursive.

By the previous proposition, the property $\fn{Deriv}(x)$ which holds
iff $x$ is the code of a correct !!{derivation}~$\pi$ in $\Log{LK}$ is
primitive recursive.  If $x$ is such a code, then $\fn{EndSequent}(x)$
is the code of the end-sequent of~$\pi$, and so
$(\fn{EndSequent}(x))_0$ is the code of the left side of the end
sequent and $(\fn{EndSequent}(x))_1$ the right side.  So we can
express ``the right side of the end-sequent of~$\pi$ is~$!A$'' as
$\len{(\fn{EndSequent}(x))_1} = 1 \land ((\fn{EndSequent}(x))_1)_0 =
x$.  The left side of the end-sequent of $\pi$ is of course
automatically finite, we just have to express that every sentence in
it is in~$\Gamma$.  Thus we can define $\Prf[\Gamma](x, y)$ by
\begin{align*}
\Prf[\Gamma](x, y) \defiff {}&
\fn{Deriv}(x) \land {} \\
& \bforall{i <
  \len{(\fn{EndSequent}(x))_0}}{R_\Gamma(((\fn{EndSequent}(x))_0)_i)} \land {}\\
& \len{(\fn{EndSequent}(x))_1} = 1 \land ((\fn{EndSequent}(x))_1)_0 = y.
\end{align*}
\end{proof}


% ==== OLP file: content/incompleteness/arithmetization-syntax/proofs-in-nd.tex ====
% Part: incompleteness
% Chapter: arithmetization-syntax
% Section: proofs-in-nd



\olfileid{inc}{art}{pnd}
\olsection{\usetoken{P}{derivation} in Natural Deduction}

\begin{explain}
In order to arithmetize !!{derivation}s, we must represent
!!{derivation}s as numbers. Since !!{derivation}s are trees of
!!{formula}s where each inference carries one or two labels, a
recursive representation is the most obvious approach: we represent a
!!{derivation} as a tuple, the components of which are the number of
immediate sub-!!{derivation}s leading to the premises of the last
inference, the representations of these sub-!!{derivation}s, and the
end-!!{formula}, the discharge label of the last inference, and a
number indicating the type of the last inference.
\end{explain}

\begin{defn}
  If $\delta$ is !!a{derivation} in natural deduction, then $\Gn{\delta}$ is
  defined inductively as follows:
  \begin{enumerate}
  \item 
    If $\delta$ consists only of the assumption~$!A$, then $\Gn{\delta}$
    is $\tuple{0, \Gn{!A}, n}$. The number~$n$ is~$0$ if it is an
    !!{undischarged} assumption, and the numerical label otherwise.
  \item 
    If $\delta$ ends in an inference with zero, one, two, or three premises,
    then $\Gn{\delta}$ is
    \begin{align*}
      & \tuple{0, \Gn{!A}, n, k}, \\
      & \tuple{1, \Gn{\delta_1}, \Gn{!A}, n, k}, \\
      & \tuple{2, \Gn{\delta_1}, \Gn{\delta_2}, \Gn{!A}, n, k}, \text{ or}\\
      & \tuple{3, \Gn{\delta_1}, \Gn{\delta_2}, \Gn{\delta_3}, \Gn{!A}, n,
        k},
    \end{align*}
    respectively.  Here $\delta_1$, $\delta_2$, $\delta_3$ are the
    sub-!!{derivation}s ending in the premise(s) of the last inference
    in $\delta$, $!A$ is the conclusion of the last inference
    in~$\delta$, $n$ is the discharge label of the last inference ($0$
    if the inference does not discharge any assumptions), and $k$~is
    given by the following table according to which rule was used in the
    last inference.
  
    \begin{tabular}{lccccccc}
      \text{Rule:} & \Intro{\land} & \Elim{\land} & \Intro{\lor} & \Elim{\lor} \\
      $k$: & 1 & 2 & 3 & 4  \\[1.5ex]
      \text{Rule:} &  \Intro{\lif} & \Elim{\lif} & \Intro{\lnot} & \Elim{\lnot} \\
      $k$: & 5 & 6 & 7 & 8 \\[1.5ex]
      \text{Rule:}  &  \FalseInt & \FalseCl & \Intro{\lforall} & \Elim{\lforall} \\
      $k$: & 9 & 10 & 11 & 12 \\[1.5ex]
      \text{Rule:} & \Intro{\lexists} & \Elim{\lexists} & \Intro{\eq} & \Elim{\eq} \\
      $k$: & 13 & 14 & 15 & 16 
    \end{tabular}
  \end{enumerate}
\end{defn}

\begin{ex}
  Consider the very simple !!{derivation}
  \begin{prooftree}
    \AxiomC{$\Discharge{!A \land !B}{1}$}
    \RightLabel{\Elim{\land}}
    \UnaryInfC{$!A$}
    \DischargeRule{\Intro{\lif}}{1}
    \UnaryInfC{$(!A \land !B) \lif !A$}
  \end{prooftree}
  The G\"odel number of the assumption would be $d_0 = \tuple{0,
    \Gn{!A \land !B}, 1}$.  The G\"odel number of the !!{derivation}
  ending in the conclusion of~$\Elim{\land}$ would be $d_1 = \tuple{1,
     d_0, \Gn{!A}, 0, 2}$ ($1$ since $\Elim{\land}$ has one premise,
  the G\"odel number of conclusion~$!A$, $0$ because no assumption is
  discharged, and $2$ is the number coding $\Elim{\land}$). The G\"odel
  number of the entire !!{derivation} then is
  $\tuple{1, d_1, \Gn{((!A \land !B) \lif
      !A)}, 1, 5}$, i.e.,
  \[
  \tuple{1, \tuple{1, \tuple{0, \Gn{(!A \land !B)}, 1}, \Gn{!A}, 0, 2},
    \Gn{((!A \land !B) \lif !A)}, 1, 5}.
  \]
\end{ex}

\begin{explain}
Having settled on a representation of !!{derivation}s, we must also
show that we can manipulate G\"odel numbers of such !!{derivation}s
primitive recursively, and express their essential properties and
relations.  Some operations are simple: e.g., given a G\"odel
number~$d$ of !!a{derivation}, $\fn{EndFmla}(d) = (d)_{(d)_0+1}$
gives us the G\"odel number of its end-!!{formula},
$\fn{DischargeLabel}(d) = (d)_{(d)_0 + 2}$ gives us the discharge
label and $\fn{LastRule}(d) = (d)_{(d)_0 + 3}$ the number indicating
the type of the last inference.  Some are much harder.  We'll at least
sketch how to do this.  The goal is to show that the relation
``$\delta$ is !!a{derivation} of~$!A$ from~$\Gamma$'' is a primitive
recursive relation of the G\"odel numbers of $\delta$ and~$!A$.
\end{explain}

\begin{prop}
  The following relations are primitive recursive:
  \begin{enumerate}
  \item $!A$ occurs as an assumption in~$\delta$ with label~$n$.
  \item All assumptions in $\delta$ with label~$n$ are of the form~$!A$
    (i.e., we can !!{discharge} the assumption~$!A$ using label~$n$
    in~$\delta$).
  \end{enumerate}
\end{prop}

\begin{proof}
  We have to show that the corresponding relations between G\"odel
  numbers of !!{formula}s and G\"odel numbers of !!{derivation}s are
  primitive recursive.
  \begin{enumerate}
  \item We want to show that $\fn{Assum}(x, d, n)$, which holds if $x$
    is the G\"odel number of an assumption of the !!{derivation} with
    G\"odel number~$d$ labelled~$n$, is primitive recursive.  This is
    the case if the !!{derivation} with G\"odel number~$\tuple{0, x,
      n}$ is a sub-!!{derivation} of~$d$. Note that the way we code
    !!{derivation}s is a special case of the coding of trees introduced in
    \olref[cmp][rec][tre]{sec}, so the primitive recursive function
    $\fn{SubtreeSeq}(d)$ gives a sequence of G\"odel numbers of all
    sub-!!{derivation}s of~$d$ (of length at most $d$). So we can
    define
    \[
    \fn{Assum}(x, d, n) \defiff \bexists{i<d}{(\fn{SubtreeSeq}(d))_i =
      \tuple{0, x, n}}.
    \]
  \item We want to show that $\fn{Discharge}(x, d, n)$, which holds if
    all assumptions with label~$n$ in the !!{derivation} with G\"odel
    number~$d$ all are the !!{formula} with G\"odel number~$x$.  But
    this relation holds iff $\bforall{y<d}{(\fn{Assum}(y, d, n) \lif y
      = x)}$.
  \end{enumerate}
\end{proof}

\begin{prop}
  \ollabel{prop:followsby}
  The property $\fn{Correct}(d)$ which holds iff the last inference in
  the !!{derivation}~$\delta$ with G\"odel number~$d$ is correct, is
  primitive recursive.
\end{prop}

\begin{proof}
  Here we have to show that for each rule of inference~$R$ the
  relation $\fn{FollowsBy}_R(d)$ is primitive recursive, where
  $\fn{FollowsBy}_R(d)$ holds iff $d$ is the G\"odel number of
  !!{derivation}~$\delta$, and the end-!!{formula} of~$\delta$ follows
  by a correct application of~$R$ from the immediate
  sub-!!{derivation}s of~$\delta$.

  A simple case is that of the \Intro{\land} rule. If $\delta$ ends in
  a correct $\Intro{\land}$ inference, it looks like this:
  \begin{prooftree}
    \AxiomC{}
    \RightLabel{$\delta_1$}
    \DeduceC{$!A$}
    
    \AxiomC{}
    \RightLabel{$\delta_2$}
    \DeduceC{$!B$}

    \RightLabel{\Intro\land}
    \BinaryInfC{$!A \land !B$}
  \end{prooftree}
  Then the G\"odel number~$d$ of~$\delta$ is $\tuple{2, d_1, d_2,
    \Gn{(!A \land !B)}, 0, k}$ where $\fn{EndFmla}(d_1) = \Gn{!A}$,
  $\fn{EndFmla}(d_2) = \Gn{!B}$, $n=0$, and $k=1$. So we can define
  $\fn{FollowsBy}_{\Intro\land}(d)$ as
  \begin{multline*}
    (d)_0 = 2 \land \fn{DischargeLabel}(d) = 0 \land \fn{LastRule}(d) = 1 \land {}\\
    \fn{EndFmla}(d) = {}\\
    \Gn{(} \concat \fn{EndFmla}((d)_1) \concat \Gn{\land}
    \concat \fn{EndFmla}((d)_2) \concat \Gn{)}.
  \end{multline*}

  Another simple example is the $\Intro\eq$ rule. This has no premises,
  so $(d)_0 = 0$, like assumptions. It also has no discharge label,
  i.e., $n=0$.  However, $!A$ must be of the form $\eq[t][t]$, for a
  closed term~$t$. Here, a primitive recursive definition is
  \begin{multline*}
    (d)_0 = 0 \land \fn{DischargeLabel}(d) = 0 \land {}\\
    \bexists{t<d}{(\fn{ClTerm}(t) \land 
      \fn{EndFmla}(d) = {}\\
      \Gn{{\eq}(} \concat t \concat \Gn{,} \concat t \concat \Gn{)})}.
  \end{multline*}

  For a more complicated example, $\fn{FollowsBy}_{\Intro{\lif}}(d)$
  holds iff the end-!!{formula} of~$\delta$ is of the form $(!A \lif
  !B)$, where the end-!!{formula} of $\delta_1$ is~$!B$, and any
  assumption in~$\delta$ labelled~$n$ is of the form~$!A$.  We can
  express this primitive recursively by
  \begin{multline*}
    (d)_0 = 1 \land {}\\
    \bexists{a<d}{(\fn{Discharge}(a, (d)_1, \fn{DischargeLabel}(d)) \land {}}\\
      \fn{EndFmla}(d) = (\Gn{(} \concat a \concat \Gn{\lif}
      \concat \fn{EndFmla}((d)_1) \concat \Gn{)}))
  \end{multline*}
  (Think of $a$ as the G\"odel number of~$!A$).

  For another example, consider \Intro{\lexists}.  Here, the last
  inference in~$\delta$ is correct iff there is !!a{formula}~$!A$, a
  closed term~$t$ and !!a{variable}~$x$ such that $\Subst{!A}{t}{x}$
  is the end-!!{formula} of the !!{derivation}~$\delta_1$ and
  $\lexists[x][!A]$ is the conclusion of the last inference.  So,
  $\fn{FollowsBy}_{\Intro{\lexists}}(d)$ holds iff
  \begin{multline*}
    (d)_0 = 1 \land \fn{DischargeLabel}(d) = 0 \land {} \\
    \bexists{a < d}{\bexists{x<d}{\bexists{t<d}{
          (\fn{ClTerm}(t) \land \fn{Var}(x)  \land {}}}}\\
    \fn{Subst}(a,t,x) = \fn{EndFmla}((d)_1) \land
    \fn{EndFmla}(d) = (\Gn{\lexists} \concat x \concat a)).
  \end{multline*}

  We then define $\fn{Correct}(d)$ as
  \begin{multline*}
    \fn{Sent}(\fn{EndFmla}(d)) \land {}\\
    (\fn{LastRule}(d) = 1 \land
    \fn{FollowsBy}_{\Intro\land}(d)) \lor \dots \lor {}\\
    (\fn{LastRule}(d) = 16 \land \fn{FollowsBy}_{\Elim\eq}(d)) \lor {}\\
    \bexists{n<d}{\bexists{x<d}{(d = \tuple{0, x, n})}}.
  \end{multline*}
  The first line ensures that the end-!!{formula} of~$d$ is a
  sentence. The last line covers the case where $d$ is just an
  assumption.
\end{proof}

\begin{prob}
  Define the following properties as in
  \olref[inc][art][pnd]{prop:followsby}:
  \begin{enumerate}
  \item $\fn{FollowsBy}_{\Elim{\lif}}(d)$,
  \item $\fn{FollowsBy}_{\Elim{\eq}}(d)$,
  \item $\fn{FollowsBy}_{\Elim{\lor}}(d)$,
  \item $\fn{FollowsBy}_{\Intro{\lforall}}(d)$.
  \end{enumerate}
  For the last one, you will have to also show that you can test
  primitive recursively if the last inference of the
  !!{derivation} with G\"odel number~$d$ satisfies the eigenvariable
  condition, i.e., the eigenvariable~$a$ of the $\Intro{\lforall}$
  inference occurs neither in the end-!!{formula} of~$d$ nor in an
  open assumption of~$d$. You may use the primitive recursive
  predicate $\fn{OpenAssum}$ from
  \olref[inc][art][pnd]{prop:openassum} for this.
\end{prob}

\begin{prop}
  \ollabel{prop:deriv}
  The relation $\fn{Deriv}(d)$ which holds if $d$ is the G\"odel
  number of a correct !!{derivation}~$\delta$, is primitive recursive.
\end{prop}

\begin{proof}
  !!^a{derivation}~$\delta$ is correct if every one of its inferences
  is a correct application of a rule, i.e., if every one of its
  sub-!!{derivation}s ends in a correct inference. So, $\fn{Deriv}(d)$
  iff
  \[
  \bforall{i<\len{\fn{SubtreeSeq}(d)}}{\fn{Correct}((\fn{SubtreeSeq}(d))_i)}
  \]
\end{proof}

\begin{prop}
  \ollabel{prop:openassum} The relation $\fn{OpenAssum}(z, d)$ that
  holds if $z$ is the G\"odel number of !!a{undischarged} assumption~$!A$
  of the !!{derivation}~$\delta$ with G\"odel number~$d$, is primitive
  recursive.
\end{prop}

\begin{proof}
  An occurrence of an assumption is !!{discharged} if it occurs with
  label~$n$ in a sub-!!{derivation} of~$\delta$ that ends in a rule
  with discharge label~$n$. So $!A$ is !!a{undischarged} assumption
  of~$\delta$ if at least one of its occurrences is not !!{discharged}
  in~$\delta$. We must be careful: $\delta$ may contain both
  !!{discharged} and !!{undischarged} occurrences of~$!A$.

  Consider a sequence $\delta_0$, \dots, $\delta_k$ where $\delta_0 =
  \delta$, $\delta_k$ is the assumption $\Discharge{!A}{n}$
  (for some~$n$), and $\delta_{i+1}$ is an immediate sub-!!{derivation}
  of~$\delta_i$. If such a sequence exists in which no $\delta_i$
  ends in an inference with discharge label~$n$, then $!A$ is
  !!a{undischarged} assumption of~$\delta$.

  The primitive recursive function $\fn{SubtreeSeq}(d)$ provides us
  with a sequence of G\"odel numbers of all sub-!!{derivation}s
  of~$\delta$. Any sequence of G\"odel numbers of sub-!!{derivation}s
  of~$\delta$ is a subsequence of it. Being a subsequence of is a
  primitive recursive relation: $\fn{Subseq}(s, s')$ holds iff
  $\bforall{i<\len{s}}{\lexists[j<\len{s'}][(s)_i = (s')_j]}$. Being an
  immediate sub-!!{derivation} is as well: $\fn{Subderiv}(d, d')$ iff
  $\bexists{j<(d')_0}{d = (d')_j}$. So we can define
  $\fn{OpenAssum}(z, d)$ by
  \begin{multline*}
    \bexists{s<\fn{SubtreeSeq}(d)}{(\fn{Subseq}(s, \fn{SubtreeSeq}(d))
      \land (s)_0 = d \land {}} \\
    \bexists{n<d}{((s)_{\len{s} \tsub 1} = \tuple{0, z, n} \land {}}\\
      \bforall{i<(\len{s} \tsub 1)}{(\fn{Subderiv}((s)_{i+1}, (s)_i) \land {}}\\
      \fn{DischargeLabel}((s)_i) \neq n))).
  \end{multline*}
\end{proof}

\begin{prop}
  \ollabel{prop:prf-prim-rec}
  Suppose $\Gamma$ is a primitive recursive set of !!{sentence}s.
  Then the relation $\Prf[\Gamma](x, y)$ expressing ``$x$ is the code
  of !!a{derivation}~$\delta$ of $!A$ from !!{undischarged}
  assumptions in~$\Gamma$ and $y$ is the G\"odel number of~$!A$'' is
  primitive recursive.
\end{prop}

\begin{proof}
  Suppose ``$y \in \Gamma$'' is given by the primitive recursive
  predicate~$R_\Gamma(y)$.  We have to show that $\Prf[\Gamma](x, y)$
  which holds iff $y$ is the G\"odel number of a sentence~$!A$ and
  $x$~is the code of a natural deduction !!{derivation} with end
  !!{formula}~$!A$ and all !!{undischarged} assumptions in~$\Gamma$ is
  primitive recursive.

  By \olref{prop:deriv}, the property $\fn{Deriv}(x)$ which holds iff
  $x$ is the G\"odel number of a correct !!{derivation}~$\delta$ in
  natural deduction is primitive recursive. Thus we can define
  $\Prf[\Gamma](x, y)$ by
  \begin{align*}
    \Prf[\Gamma](x, y) \defiff {}
    & \fn{Deriv}(x) \land \fn{EndFmla}(x) = y \land {} \\
    & \bforall{z < x}{(\fn{OpenAssum}(z, x) \lif R_\Gamma(z))}.
  \end{align*}
\end{proof}


% ==== OLP file: content/incompleteness/arithmetization-syntax/substitution.tex ====
% Part: incompleteness
% Chapter: arithmetization-syntax
% Section: substitution



\olfileid{inc}{art}{sub}
\olsection{Substitution}

Recall that substitution is the operation of replacing all free
occurrences of !!a{variable}~$u$ in !!a{formula}~$!A$ by a term~$t$,
written $\Subst{!A}{t}{u}$. This operation, when carried out on
G\"odel numbers of !!{variable}s, !!{formula}s, and terms, is
primitive recursive.

\begin{prop}\ollabel{prop:subst-primrec}
There is a primitive recursive function $\fn{Subst}(x, y, z)$ with the
property that
\[
\fn{Subst}(\Gn{!A}, \Gn{t}, \Gn{u}) = \Gn{\Subst{!A}{t}{u}}.
\]
\end{prop}

\begin{proof}
We can then define a function $\fn{hSubst}$ by primitive recursion as
follows:
\begin{multline*}
\begin{aligned}
\fn{hSubst}(x, y, z, 0) & = \emptyseq \\
\fn{hSubst}(x, y, z, i+1) & =
\end{aligned}\\
\begin{cases}
\fn{hSubst}(x, y, z, i) \concat y & \text{if $\fn{FreeOcc}(x, z, i)$} \\
\fn{append}(\fn{hSubst}(x, y, z, i), (x)_{i}) & \text{otherwise.}
\end{cases}
\end{multline*}
$\fn{Subst}(x, y, z)$ can now be defined as $\fn{hSubst}(x, y, z, \len{x})$.
\end{proof}

\begin{prop}
\ollabel{prop:free-for}
The relation $\fn{FreeFor}(x, y, z)$, which holds iff the term with
G\"odel number~$y$ is !!{free for} the variable with G\"odel
number~$z$ in the formula with G\"odel number~$x$, is primitive
recursive.
\end{prop}

\begin{proof} Exercise. \end{proof}

\begin{prob}
Prove \olref[inc][art][sub]{prop:free-for}
\end{prob}

\end{document}

\end{document}

\end{document}
